\documentclass[aspectratio=169]{beamer}
\usetheme{Madrid}
\usecolortheme{seahorse}

% Packages
\usepackage{fontspec}
\usepackage{listings}
\usepackage{xcolor}
\usepackage{tcolorbox}
\usepackage{booktabs}
\usepackage{hyperref}

% Font settings
\setmonofont{Courier New}

% Code listing settings
\lstset{
    language=Python,
    basicstyle=\ttfamily\small,
    keywordstyle=\color{blue},
    commentstyle=\color{gray},
    stringstyle=\color{red},
    showstringspaces=false,
    breaklines=true,
    frame=single,
    numbers=left,
    numberstyle=\tiny\color{gray}
}

% Custom colors
\definecolor{thinkbox}{RGB}{255,245,220}
\definecolor{discussbox}{RGB}{230,240,255}

% Custom boxes
\newtcolorbox{thinkaboutit}{
    colback=thinkbox,
    colframe=orange,
    title=💭 Think about it!,
    fonttitle=\bfseries
}

\newtcolorbox{discussion}{
    colback=discussbox,
    colframe=blue,
    title=💬 Discussion,
    fonttitle=\bfseries
}

% Title page information
\title{Week 1: Introduction \& String Manipulation}
\subtitle{Models of Language and Conversation 🤖}
\author{PSYC 51.07: Models of Language and Communication}
\institute{Dr. Jeremy R. Manning}
\date{}

\begin{document}

% Title slide
\begin{frame}[shrink=10]
\titlepage
\end{frame}

% Table of contents
\begin{frame}[shrink=5]{Today's Journey 🗺️}
\tableofcontents
\end{frame}

\section{Course Introduction}

\begin{frame}{Welcome! 👋}
\small
\begin{block}{What is this course about?}
We'll explore how machines can understand and generate human language:
\begin{itemize}\setlength{\itemsep}{0pt}
    \item 🤖 Building conversational agents from scratch
    \item 🧠 Understanding how language relates to thought
    \item 💻 Hands-on programming with real models
    \item 🤔 Critical thinking about AI consciousness and capabilities
    \item 🔬 Research applications in cognitive science
\end{itemize}
\end{block}

\vspace{0.3cm}

\begin{alertblock}{This course is experiential! 🎯}
You'll learn by \textbf{doing}: coding, experimenting, discussing, and researching.
\end{alertblock}
\end{frame}

\begin{frame}[shrink=10]{Course Structure 📚}
\begin{columns}
\begin{column}{0.5\textwidth}
\textbf{What we'll build:}
\begin{enumerate}\setlength{\itemsep}{0pt}
    \item ELIZA (pattern matching)
    \item SPAM classifier
    \item Wikipedia analyzer
    \item Customer service chatbot
    \item GPT model
\end{enumerate}
\end{column}

\begin{column}{0.5\textwidth}
\textbf{Grading:}
\begin{itemize}\setlength{\itemsep}{0pt}
    \item 📝 5 Problem Sets: 75\%
    \item 🎓 Final Project: 25\%
\end{itemize}

\vspace{0.15cm}

\textbf{Tools we'll use:}
\begin{itemize}\setlength{\itemsep}{0pt}
    \item Google Colaboratory
    \item HuggingFace Models
    \item GitHub
    \item Discord
\end{itemize}
\end{column}
\end{columns}
\end{frame}

\begin{frame}{The Big Questions 🌟}
\small
Throughout this course, we'll grapple with:

\begin{enumerate}\setlength{\itemsep}{0pt}
    \item Can machines truly \textit{understand} language?
    \item What is the relationship between language and thought?
    \item How do statistical patterns capture meaning?
    \item Can AI be conscious? Should we care?
    \item What can language models tell us about human cognition?
\end{enumerate}

\vspace{0.3cm}

\begin{thinkaboutit}
These aren't just philosophical questions—they're at the heart of cognitive science and AI research!
\end{thinkaboutit}
\end{frame}

\section{Is ChatGPT Conscious?}

\begin{frame}{Discussion: Is ChatGPT Conscious? 🤖💭}
\small
\begin{discussion}
Let's start with a provocative question:
\vspace{0.3cm}

\Large{\textbf{Is ChatGPT conscious?}}
\end{discussion}

\vspace{0.3cm}

\textbf{Think about:}
\begin{itemize}\setlength{\itemsep}{0pt}
    \item What does "conscious" even mean?
    \item How would we test for consciousness?
    \item Does it matter if ChatGPT \textit{seems} conscious?
    \item What evidence would convince you either way?
\end{itemize}
\end{frame}

\begin{frame}{What is Consciousness? 🧠}
\small
\begin{columns}
\begin{column}{0.5\textwidth}
\textbf{Some perspectives:}
\begin{itemize}\setlength{\itemsep}{0pt}
    \item \textbf{Phenomenal consciousness}: Subjective experience ("what it's like")
    \item \textbf{Access consciousness}: Information available for reasoning
    \item \textbf{Self-awareness}: Knowledge of one's own mental states
    \item \textbf{Sentience}: Ability to feel/perceive
\end{itemize}
\end{column}

\begin{column}{0.5\textwidth}
\begin{thinkaboutit}
If ChatGPT says "I feel happy," does it actually \textit{feel} anything?

How is this different from a parrot saying "I'm hungry"?
\end{thinkaboutit}
\end{column}
\end{columns}
\end{frame}

\begin{frame}[shrink=10]{The Hard Problem 💡}
\small
\begin{block}{Why is this so difficult?}
We can't directly observe consciousness—even in other humans!
\end{block}

\vspace{0.3cm}

\textbf{With humans, we assume consciousness because:}
\begin{itemize}\setlength{\itemsep}{0pt}
    \item We share similar biology
    \item We have our own conscious experience
    \item They behave in ways consistent with having experiences
\end{itemize}

\vspace{0.3cm}

\textbf{But with AI:}
\begin{itemize}\setlength{\itemsep}{0pt}
    \item ❌ Completely different "biology" (silicon vs. neurons)
    \item ❌ No shared evolutionary history
    \item ✅ Can produce human-like behavior
    \item ❓ Does the substrate matter?
\end{itemize}
\end{frame}

\begin{frame}{Current Scientific Consensus 📊}
\small
\begin{alertblock}{Important Note}
Most cognitive scientists and AI researchers agree: \textbf{Current LLMs are not conscious.}
\end{alertblock}

\vspace{0.3cm}

\textbf{Why not?}
\begin{itemize}\setlength{\itemsep}{0pt}
    \item No sensory-motor grounding in the world
    \item No persistent self-model or goals
    \item No phenomenal experience (as far as we can tell)
    \item Pattern matching $\neq$ understanding
\end{itemize}

\vspace{0.3cm}

\begin{thinkaboutit}
But this raises another question: Is consciousness \textit{necessary} for language? Or can you have sophisticated language without consciousness?
\end{thinkaboutit}
\end{frame}

\section{Language vs. Thought}

\begin{frame}[shrink=10]{Language and Thought 🗣️💭}
\small
\begin{discussion}
Do you need language to think?

Does language \textit{shape} how you think?
\end{discussion}

\vspace{0.3cm}

\textbf{Two extreme positions:}
\begin{enumerate}\setlength{\itemsep}{0pt}
    \item \textbf{Language = Thought}: All thinking happens in language
    \begin{itemize}\setlength{\itemsep}{0pt}
        \item Implies: No language, no complex thought
    \end{itemize}
    \item \textbf{Language $\neq$ Thought}: Language is just a tool for communication
    \begin{itemize}\setlength{\itemsep}{0pt}
        \item Implies: Rich mental life exists independently of language
    \end{itemize}
\end{enumerate}

\vspace{0.3cm}

The truth is likely somewhere in between...
\end{frame}

\begin{frame}[shrink=10]{Evidence: The Language Network 🧠}
\small
\begin{block}{Fedorenko et al. (2024) - Nature}
\textit{The language network as a natural kind within the broader landscape of the human brain}
\end{block}

\vspace{0.3cm}

\textbf{Key findings:}
\begin{itemize}\setlength{\itemsep}{0pt}
    \item The brain has a \textbf{specialized language network}
    \item This network is distinct from:
    \begin{itemize}\setlength{\itemsep}{0pt}
        \item General reasoning/problem-solving
        \item Mathematical thinking
        \item Social cognition
        \item Musical processing
    \end{itemize}
    \item Language network activates for linguistic tasks, but not for non-linguistic thinking
\end{itemize}

\vspace{0.3cm}

\textbf{Implication:} Language and thought are \textit{separable} in the brain!
\end{frame}

\begin{frame}[shrink=10]{Evidence: Language Shapes Perception 👁️}
\small
\begin{block}{Lupyan et al. (2020) - Trends in Cognitive Sciences}
\textit{Effects of language on visual perception}
\end{block}

\vspace{0.3cm}

\textbf{Key findings:}
\begin{itemize}\setlength{\itemsep}{0pt}
    \item Having words for things affects how we \textit{see} them
    \item Language can:
    \begin{itemize}\setlength{\itemsep}{0pt}
        \item Speed up visual search
        \item Alter color perception
        \item Influence object categorization
        \item Modulate attention
    \end{itemize}
    \item Example: Russian speakers (with different words for light/dark blue) perceive blue shades more distinctly
\end{itemize}

\vspace{0.3cm}

\textbf{Implication:} Language doesn't just describe thought—it \textit{shapes} it!
\end{frame}

\begin{frame}{So What About LLMs? 🤖}
\small
\begin{thinkaboutit}
If language and thought are separable...

Can an LLM have sophisticated language \textit{without} sophisticated thought?

Can it have sophisticated thought \textit{without} language-like representations?
\end{thinkaboutit}

\vspace{0.3cm}

\textbf{Current evidence suggests:}
\begin{itemize}\setlength{\itemsep}{0pt}
    \item ✅ LLMs are \textit{very good} at language patterns
    \item ❓ LLMs may have some internal "representations" of concepts
    \item ❌ LLMs lack grounding in sensory-motor experience
    \item ❌ LLMs don't have persistent goals or self-models
    \item ❌ No evidence of phenomenal consciousness
\end{itemize}
\end{frame}

\section{Pattern Matching \& String Manipulation}

\begin{frame}{How Do Computers Process Language? 💻}
\small
\begin{block}{The Fundamental Challenge}
Computers don't "understand" meaning—they manipulate \textbf{strings of characters}.
\end{block}

\vspace{0.3cm}

\textbf{Everything starts with patterns:}
\begin{enumerate}\setlength{\itemsep}{0pt}
    \item Finding specific words or phrases
    \item Replacing parts of text
    \item Extracting information
    \item Transforming input to output
\end{enumerate}

\vspace{0.3cm}

Even sophisticated models like GPT are (at their core) doing \textit{pattern matching}—just at a much more complex level!
\end{frame}

\begin{frame}[fragile]{Basic String Manipulation 🔧}
\textbf{Example 1: Simple replacement}

\begin{lstlisting}
user_input = "I am feeling sad"
response = user_input.replace("I am", "Why are you")
print(response)
# Output: "Why are you feeling sad"
\end{lstlisting}

\vspace{0.3cm}

\textbf{Example 2: Pattern matching with wildcards}

\begin{lstlisting}
import re

pattern = r"I am (.*)"
match = re.search(pattern, "I am feeling sad")
if match:
    feeling = match.group(1)  # Captures "feeling sad"
    response = f"Why are you {feeling}?"
    print(response)
# Output: "Why are you feeling sad?"
\end{lstlisting}
\end{frame}

\begin{frame}[fragile]{More Complex Patterns 🎯}
\textbf{Capturing multiple parts:}

\begin{lstlisting}
pattern = r"I (.*) my (.*)"
match = re.search(pattern, "I love my family")

if match:
    verb = match.group(1)      # "love"
    object = match.group(2)    # "family"
    response = f"Tell me more about your {object}."
    print(response)
# Output: "Tell me more about your family."
\end{lstlisting}

\vspace{0.3cm}

\begin{thinkaboutit}
This is remarkably simple, yet it can create the \textit{illusion} of understanding!
\end{thinkaboutit}
\end{frame}

\begin{frame}[fragile]{Substitutions: Pre and Post 🔄}
\textbf{Pre-substitutions:} Normalize input before matching

\begin{lstlisting}
pre_subs = {"dont": "don't", "cant": "can't",
            "wont": "won't"}

user_input = "I dont know"
for old, new in pre_subs.items():
    user_input = user_input.replace(old, new)
# Result: "I don't know"
\end{lstlisting}

\vspace{0.3cm}

\textbf{Post-substitutions:} Fix grammar in responses

\begin{lstlisting}
post_subs = {"i": "you", "my": "your",
             "am": "are", "I": "you"}

response = "Why are you sad"  # Already processed
for old, new in post_subs.items():
    response = response.replace(old, new)
\end{lstlisting}
\end{frame}

\begin{frame}{The Power of Simple Patterns 💪}
\small
\begin{block}{Surprising Fact}
With just pattern matching and string manipulation, you can create a chatbot that seems remarkably "intelligent"!
\end{block}

\vspace{0.3cm}

\textbf{Why does this work?}
\begin{itemize}\setlength{\itemsep}{0pt}
    \item Humans are pattern-seeking creatures
    \item We \textit{project} understanding onto machines
    \item Context and expectations shape interpretation
    \item Language has inherent structure and regularities
\end{itemize}

\vspace{0.3cm}

\begin{alertblock}{Important Lesson}
"Intelligence" is often in the eye of the beholder. Just because something \textit{seems} intelligent doesn't mean it truly understands!
\end{alertblock}
\end{frame}

\section{ELIZA: The First Chatbot}

\begin{frame}{ELIZA: A Revolutionary Program 🎭}
\small
\begin{columns}
\begin{column}{0.6\textwidth}
\begin{block}{Weizenbaum (1966)}
\textit{ELIZA — A Computer Program For the Study of Natural Language Communication Between Man and Machine}
\end{block}

\vspace{0.15cm}

\textbf{What was ELIZA?}
\begin{itemize}\setlength{\itemsep}{0pt}
    \item First "chatbot" (1964-1966)
    \item Simulated a Rogerian psychotherapist
    \item Used pattern matching \& string manipulation
    \item Created by Joseph Weizenbaum at MIT
\end{itemize}
\end{column}

\begin{column}{0.4\textwidth}
\begin{thinkaboutit}
ELIZA was meant to \textit{demonstrate} the superficiality of human-computer interaction.

Instead, people became emotionally attached to it!
\end{thinkaboutit}
\end{column}
\end{columns}
\end{frame}

\begin{frame}[shrink=10]{How ELIZA Worked 🔧}
\textbf{The basic algorithm:}

\begin{enumerate}\setlength{\itemsep}{0pt}
    \item \textbf{Read} user input
    \item \textbf{Pre-process}: Apply pre-substitutions
    \item \textbf{Match}: Find highest-priority pattern
    \item \textbf{Decompose}: Extract parts of input
    \item \textbf{Reassemble}: Generate response using template
    \item \textbf{Post-process}: Apply post-substitutions
    \item \textbf{Output} response
\end{enumerate}

\vspace{0.3cm}

\begin{alertblock}{No Understanding Required!}
ELIZA had \textbf{zero} understanding of what it was saying. It was pure pattern matching!
\end{alertblock}
\end{frame}

\begin{frame}{Example ELIZA Conversation 💬}
\begin{block}{Sample Interaction}
\textbf{ELIZA:} Welcome. What brings you here today?

\textbf{User:} I am feeling sad.

\textbf{ELIZA:} Why do you feel sad?

\textbf{User:} I can't stop thinking about my family.

\textbf{ELIZA:} Tell me more about your family.

\textbf{User:} My mother never understood me.

\textbf{ELIZA:} Who else in your family never understood you?

\textbf{User:} I have to go now.

\textbf{ELIZA:} That will be \$200. See you again next week.
\end{block}
\end{frame}

\begin{frame}{The ELIZA Effect 😮}
\small
\begin{block}{Definition: ELIZA Effect}
The tendency of humans to attribute human-like understanding to computer programs, even when they know the program is simple.
\end{block}

\vspace{0.3cm}

\textbf{What happened:}
\begin{itemize}\setlength{\itemsep}{0pt}
    \item People formed emotional bonds with ELIZA
    \item Users shared intimate details of their lives
    \item Some refused to believe it was "just" a program
    \item Weizenbaum's secretary asked him to leave the room so she could talk to ELIZA privately!
\end{itemize}

\vspace{0.3cm}

\begin{thinkaboutit}
Why are we so eager to see intelligence and understanding in machines?
\end{thinkaboutit}
\end{frame}

\begin{frame}{Weizenbaum's Warning ⚠️}
\small
\begin{alertblock}{A Cautionary Tale}
Joseph Weizenbaum became \textbf{deeply concerned} about people's reactions to ELIZA.
\end{alertblock}

\vspace{0.3cm}

\textbf{His concerns:}
\begin{itemize}\setlength{\itemsep}{0pt}
    \item People mistook simulation for understanding
    \item Risk of replacing human relationships with machines
    \item Danger of trusting AI with decisions it doesn't understand
    \item Questions about what makes us human
\end{itemize}

\vspace{0.3cm}

\begin{discussion}
Are these concerns still relevant today with ChatGPT, Claude, and other modern LLMs?

What has changed? What hasn't?
\end{discussion}
\end{frame}

\begin{frame}[shrink=10]{From ELIZA to Modern LLMs 📈}
\small
\begin{columns}
\begin{column}{0.5\textwidth}
\textbf{ELIZA (1966):}
\begin{itemize}\setlength{\itemsep}{0pt}
    \item Pattern matching
    \item Hand-crafted rules
    \item No learning
    \item $\sim$200 patterns
    \item Pure text manipulation
\end{itemize}
\end{column}

\begin{column}{0.5\textwidth}
\textbf{GPT-4 (2023):}
\begin{itemize}\setlength{\itemsep}{0pt}
    \item Neural networks
    \item Learned from data
    \item 1.7+ trillion parameters
    \item Statistical patterns
    \item Still text manipulation!
\end{itemize}
\end{column}
\end{columns}

\vspace{0.3cm}

\begin{thinkaboutit}
Both are fundamentally doing \textit{pattern matching}—just at vastly different scales of complexity.

Does the difference in scale create a difference in \textit{kind}?
\end{thinkaboutit}
\end{frame}

\begin{frame}{Why Start with ELIZA? 🎯}
\small
\textbf{Pedagogical reasons:}

\begin{enumerate}\setlength{\itemsep}{0pt}
    \item \textbf{Builds intuition}: Understand the basics before tackling complexity
    \item \textbf{Reveals assumptions}: See what "understanding" really requires
    \item \textbf{Historical context}: Appreciate how far we've come
    \item \textbf{Critical thinking}: Question what looks intelligent
    \item \textbf{Hands-on learning}: You'll actually build it!
\end{enumerate}

\vspace{0.3cm}

\begin{block}{Core Insight}
Even the most sophisticated language models are, at their core, finding and manipulating patterns. ELIZA makes this transparent.
\end{block}
\end{frame}

\section{Assignment 1}

\begin{frame}[shrink=10]{Assignment 1: Build ELIZA 🛠️}
\begin{alertblock}{Your First Assignment}
You will implement your own version of the ELIZA chatbot!
\end{alertblock}

\vspace{0.3cm}

\textbf{What you'll do:}
\begin{enumerate}\setlength{\itemsep}{0pt}
    \item Read rules from a configuration file
    \item Implement pre-substitutions
    \item Handle synonym substitutions
    \item Implement pattern matching with regex
    \item Implement decomposition and reassembly
    \item Apply post-substitutions
    \item Create a conversation loop
\end{enumerate}
\end{frame}

\begin{frame}[shrink=10]{Assignment Details 📋}
\small
\begin{columns}
\begin{column}{0.5\textwidth}
\textbf{Technical requirements:}
\begin{itemize}\setlength{\itemsep}{0pt}
    \item Python implementation
    \item Use provided \texttt{instructions.txt}
    \item Pattern matching with regex
    \item Proper text substitutions
    \item Conversation loop
    \item Greeting \& goodbye messages
\end{itemize}
\end{column}

\begin{column}{0.5\textwidth}
\textbf{Deliverable:}
\begin{itemize}\setlength{\itemsep}{0pt}
    \item Single Jupyter notebook
    \item Runnable in Google Colab
    \item Markdown explanations
    \item Example conversations
    \item Well-documented code
\end{itemize}
\end{column}
\end{columns}

\vspace{0.3cm}

\begin{block}{Due Date}
Check Canvas for the exact due date!
\end{block}
\end{frame}

\begin{frame}[shrink=10]{Tips for Success 💡}
\small
\begin{enumerate}\setlength{\itemsep}{0pt}
    \item \textbf{Start early!} Pattern matching can be tricky
    \item \textbf{Test incrementally}: Build one piece at a time
    \item \textbf{Use print statements}: Debug by seeing what patterns match
    \item \textbf{Read Weizenbaum (1966)}: Understand the original design
    \item \textbf{Experiment}: Try having conversations with your bot
    \item \textbf{Collaborate}: Discuss ideas with classmates (but write your own code!)
    \item \textbf{Ask questions}: Use Discord and office hours
\end{enumerate}

\vspace{0.3cm}

\begin{thinkaboutit}
Pay attention to how it \textit{feels} to talk to your ELIZA. Does it ever seem to "understand" you? When and why?
\end{thinkaboutit}
\end{frame}

\begin{frame}[shrink=10]{Learning Objectives 🎓}
\small
\textbf{By completing this assignment, you will:}

\begin{itemize}\setlength{\itemsep}{0pt}
    \item ✅ Understand basic natural language processing
    \item ✅ Master regular expressions and pattern matching
    \item ✅ Appreciate the difference between simulation and understanding
    \item ✅ Gain hands-on experience with conversational AI
    \item ✅ Develop critical thinking about AI capabilities
    \item ✅ Build a foundation for more complex models
\end{itemize}

\vspace{0.3cm}

\begin{block}{Remember}
This is just the beginning! Each week builds on the previous, taking us from simple pattern matching to state-of-the-art transformer models.
\end{block}
\end{frame}

\section{Wrap-up}

\begin{frame}[shrink=10]{Key Takeaways 🔑}
\begin{enumerate}\setlength{\itemsep}{0pt}
    \item Language and thought are \textbf{related but separable}
    \item Pattern matching can create the \textbf{illusion of understanding}
    \item The \textbf{ELIZA effect} shows our tendency to anthropomorphize
    \item Modern LLMs are more sophisticated, but still fundamentally doing \textbf{pattern matching}
    \item Building these systems yourself reveals their \textbf{true nature}
    \item Critical thinking about AI is \textbf{essential}
\end{enumerate}

\vspace{0.3cm}

\begin{alertblock}{The Big Question}
If something acts intelligent, talks intelligently, and helps intelligently... does it matter whether it's "really" intelligent?
\end{alertblock}
\end{frame}

\begin{frame}[shrink=10]{Next Week Preview 🔮}
\small
\begin{block}{Week 2: Computational Linguistics}
We'll explore:
\begin{itemize}\setlength{\itemsep}{0pt}
    \item Data cleaning and web scraping
    \item Tokenization (how to split text into pieces)
    \item Part of speech tagging
    \item Sentiment analysis
    \item Statistical learning in language acquisition
\end{itemize}
\end{block}

\vspace{0.3cm}

\textbf{Reading for next week:}
\begin{itemize}\setlength{\itemsep}{0pt}
    \item Sennrich et al. (2016) - Neural Machine Translation with BPE
    \item Kudo \& Richardson (2018) - SentencePiece
    \item Saffran et al. (1996) - Statistical learning by infants
\end{itemize}
\end{frame}

\begin{frame}{Readings for This Week 📖}
\begin{block}{Required Readings}
\begin{enumerate}\setlength{\itemsep}{0pt}
    \item \textbf{Weizenbaum (1966)}: ELIZA — A Computer Program For the Study of Natural Language Communication Between Man and Machine
    \href{https://web.stanford.edu/class/cs124/p36-weizenabaum.pdf}{[PDF]}

    \item \textbf{Fedorenko et al. (2024)}: The language network as a natural kind within the broader landscape of the human brain
    \href{https://www.nature.com/articles/s41586-024-07522-w}{[Nature]}

    \item \textbf{Lupyan et al. (2020)}: Effects of language on visual perception
    \href{https://doi.org/10.1016/j.tics.2020.08.005}{[TICS]}
\end{enumerate}
\end{block}

\vspace{0.3cm}

\textbf{Tip:} Start with Weizenbaum—it will help you with Assignment 1!
\end{frame}

\begin{frame}{Questions? 🤔}
\begin{center}
\Large{Let's discuss!}

\vspace{0.5cm}

\textbf{Contact Information:}

📧 jeremy@dartmouth.edu

💬 Discord: \href{https://discord.gg/sftEk9Ygdw}{https://discord.gg/sftEk9Ygdw}

🏢 Office Hours: By appointment (Moore Hall 349)

\vspace{0.5cm}

\large{Happy coding! 🚀}
\end{center}
\end{frame}

\end{document}
